\section{Introduction}

Graph neural networks (GNNs) learn graph representations by repeatedly aggregating information from local neighborhoods \cite{gori2005new,kipf2017gcn,hamilton2017inductive,velivckovic2017graph,xu2019gin}. This message-passing view has become a standard tool for graph classification because it can combine node attributes, local neighborhoods, and higher-order graph structure in a single differentiable model \cite{gilmer2017neural,duvenaud2015convolutional,wu2020comprehensive,zhou2020graph}. In molecular and protein-oriented benchmarks, the final prediction may depend on both local biochemical motifs and global graph organization, so the model must preserve useful intermediate evidence until graph-level pooling \cite{ying2018hierarchical,lee2019self,cangea2018towards}. Increasing model depth or adding parallel branches can expand capacity, but it can also introduce over-smoothing, over-squashing, redundant representations, unstable optimization, and unnecessary computation \cite{li2018deeper,oono2020simple,chen2023revisiting,alon2020bottleneck,topping2021understanding,zhao2020pairnorm}.

Residual connections are a common response to these difficulties. They improve gradient flow and make it possible to reuse earlier representations, a principle established in deep convolutional networks and later adapted to graph architectures \cite{he2016deep,bresson2017residual,li2019deepgcns,li2021training}. Other depth-control strategies, including normalization, DropEdge, APPNP-style propagation, Jumping Knowledge aggregation, and graph-specific normalization, also aim to preserve useful information across propagation steps \cite{rong2019dropedge,klicpera2018combining,xu2018representation,cai2021graphnorm,chenSimpleDeepGraph2020}. However, residual connectivity is still commonly treated as a default layer-wise component rather than as the object of study. Most residual graph classifiers reuse information along depth in the same stream. Once a model contains multiple branches, this view is incomplete: residual information may also move across branches at the same depth or through diagonal branch-depth neighborhoods.

Existing graph-classification work has explored powerful message-passing operators, hierarchical pooling, attention, and higher-order neighborhood mixing \cite{abu2019mixhop,bianchi2020graph,gao2019graph,ying2018hierarchical,beaini2020directional,dwivedi2020generalization}. These designs improve the expressive or aggregation capacity of GNNs, but they do not isolate the residual-routing question: if several branch states are available at a given depth, should the model reuse only the previous state of the same branch, exchange states laterally, or combine both directions in a local two-dimensional residual neighborhood? Answering this question requires a controlled family in which the readout, optimizer, dataset split, and message-passing backbone are held fixed while only the residual topology changes.

This paper studies residual reuse as a structured design problem on a branch-by-layer grid. The central question is not whether skip connections should be present, but where residual signals should flow. A grid view exposes three controlled neighborhoods: vertical reuse along depth, horizontal exchange across neighboring branches, and matrix-style reuse that combines same-branch, neighboring-branch, and diagonal residual paths. A sparse/gated matrix variant further tests whether matrix residual traffic should be filtered before injection. This framing follows the spirit of controlled benchmark studies \cite{keriven2020benchmark,morris2020tudataset}: the goal is not to add a large number of unrelated architectural components, but to make a single design axis explicit enough to evaluate.

The study is deliberately controlled. It fixes the graph-classification pipeline within each comparison, evaluates the same architecture family under four standard message-passing operators, and changes only the residual topology and residual-control mechanism. This separation is important because it keeps the empirical comparison focused on information-flow design rather than on unrelated changes in readout structure or task-specific feature engineering.

The work is organized around four research questions:
\begin{enumerate}
    \item \textbf{Which residual topology is strongest under a shared graph-classification protocol?} We compare plain message passing, vertical residual reuse, horizontal branch-wise exchange, local matrix residual reuse, and sparse/gated matrix reuse on six graph-classification datasets under four message-passing operators.
    \item \textbf{How does branch count change residual behavior?} We test \(B=1,\ldots,8\) for the multi-branch residual families and ask whether more branches improve accuracy monotonically or only within a useful operating region.
    \item \textbf{Why include sparse/gated matrix reuse?} We test whether matrix residual traffic should be passed densely or filtered before injection.
    \item \textbf{When does controlled matrix reuse help?} We use sensitivity scans, five-fold tuned-candidate checks, and mechanism metrics to distinguish promising residual-control settings from fold-specific artifacts.
\end{enumerate}

The main contributions are:
\begin{itemize}
    \item \textbf{A branch-by-layer residual formulation for graph classification.} Vertical, horizontal, matrix, and sparse/gated matrix residual families are expressed as local neighborhoods on the same grid, making their differences explicit while preserving a common backbone and classifier.
    \item \textbf{A controlled five-fold benchmark and branch-count ablation.} The experiments cover 24 dataset--operator combinations. MatrixRes obtains the largest number of wins, but the per-dataset winners remain mixed, and branch-count curves follow a rise-then-fall pattern.
    \item \textbf{Evidence that branch expansion is not merely a capacity knob.} Branch diversity, branch cosine similarity, and gradient norms show that additional branches help only when they create useful separation without excessive residual traffic or optimization weakening.
    \item \textbf{A conservative validation of MatrixResGated tuning.} Fold-0 sensitivity scans identify candidate operating regions, and five-fold reruns show that the smaller DD MatrixResGated setting is parameter-efficient, whereas the PROTEINS sparsity candidate is not yet a robust full-fold gain.
\end{itemize}
